\chapter{Mausereignisse \tagZusatz}
\renewcommand{\chaptertitle}{Mausereignisse}

\lehead[]{\normalfont\sffamily\hspace*{-2.00cm}\textcolor{white}{\colorbox{lightblue}{\makebox[1.60cm][r]{\thechapter}}}\hspace{0.17cm}\textcolor{lightblue}{\chaptertitle}}
\rohead[]{\textcolor{lightblue}{\chaptertitle}\normalfont\sffamily\hspace*{0.17cm}\textcolor{white}{\colorbox{lightblue}{\makebox[1.60cm][l]{\thechapter}}}\hspace{-2.00cm}}
%\chead[]{}
\rehead[]{\textcolor{lightblue}{AvHG, Inf, My}}
\lohead[]{\textcolor{lightblue}{AvHG, Inf, My}}

\lstset{style=myJava}

\section{Normale Mausereignisse}

Ebenso wie Tastaturereignisse können in einem Frame auch Mausereignisse
abgefangen werden (dazu benötigt man \lstinline|import java.awt.event.*;|). Wenn
man Mausereignisse abfangen möchte, muss man das Interface
\myClass{MouseListener} implementieren:

\medskip
\bgroup
\def\arraystretch{1.2}
\begin{tabularx}{\textwidth}{|X|p{80mm}|}
\hline
\textbf{Methoden des Interfaces \myClass{MouseListener}} & 
\textbf{Wird aufgerufen, wenn \ldots}
\\ \hline
\texttt{public void mouseClicked(MouseEvent e)} &
\ldots\ eine Maustaste gedrückt und wieder losgelassen wurde.
\\ \hline
\texttt{public void mousePressed(MouseEvent e)} &
\ldots\ eine Maustaste gedrückt wurde.
\\ \hline
\texttt{public void mouseReleased(MouseEvent e)} &
\ldots\ eine Maustaste losgelassen wurde.
\\ \hline
\texttt{public void mouseEntered(MouseEvent e)} &
\ldots\ der Mauszeiger die Komponente \glqq betreten\grqq\ hat.
\\ \hline
\texttt{public void mouseExited(MouseEvent e)} &
\ldots\ der Mauszeiger die Komponente verlassen hat.
\\ \hline
\end{tabularx}
\egroup

\medskip
Eine Klasse, die das Interface \myClass{MouseListener} implementiert, muss beim
Frame bzw.\ der Zeichenfläche durch Aufruf der folgenden Methode registriert
werden:

\begin{lstlisting}
public void addMouseListener(MouseListener l)
\end{lstlisting}

\noindent \textbf{Wichtiger Hinweis zu \myClass{HJFrame}:} Registriere den
Listener direkt auf der \lstinline|zeichenflaeche| (also
\lstinline|zeichenflaeche.addMouseListener(this);|), damit die Mauskoordinaten
exakt mit dem Koordinatensystem von \lstinline|myPaint()| übereinstimmen (ohne
Versatz durch die Titelleiste des Fensters).


\section{Mausbewegungen}

Für Mausbewegungen gibt es ein eigenes Interface \myClass{MouseMotionListener}:

\medskip
\bgroup
\def\arraystretch{1.2}
\begin{tabularx}{\textwidth}{|X|p{80mm}|}
\hline
\textbf{Methoden des Interfaces} 

\textbf{\myClass{MouseMotionListener}} & 
\textbf{Wird aufgerufen, wenn \ldots}
\\ \hline
\texttt{public void mouseDragged(MouseEvent e)} &
\ldots\ die Maus bei gedrückter Maustaste bewegt wurde.
\\ \hline
\texttt{public void mouseMoved(MouseEvent e)} &
\ldots\ die Maus ohne gedrückte Maustaste bewegt wurde.
\\ \hline
\end{tabularx}
\egroup

\medskip
Eine Klasse, die das Interface \myClass{MouseMotionListener} implementiert, muss
ebenfalls registriert werden (bei \myClass{HJFrame} wieder am besten über
\lstinline|zeichenflaeche.addMouseMotionListener(this);|):

\begin{lstlisting}
public void addMouseMotionListener(MouseMotionListener l)
\end{lstlisting}


\section{Details über Mausereignisse erfahren}

Alle diese Methoden erhalten als Parameter ein Objekt der Klasse
\myClass{MouseEvent}, in dem Details über das Mausereignis gespeichert sind.
Die Klasse \myClass{MouseEvent} besitzt unter anderem folgende Methoden:

\medskip
\bgroup
\def\arraystretch{1.2}
\begin{tabularx}{\textwidth}{|X|p{80mm}|}
\hline
\textbf{Methoden der Klasse \myClass{MouseEvent} (Auswahl)} & 
\textbf{Beschreibung}
\\ \hline
\texttt{public int getButton()} &
Gibt an, welcher Knopf gedrückt wurde:
\texttt{MouseEvent.BUTTON1} (links), \texttt{MouseEvent.BUTTON2} (Mitte) oder
\texttt{MouseEvent.BUTTON3} (rechts).
\\ \hline
\texttt{public int getClickCount()} &
Anzahl der Mausklicks (z.\,B.\ $2$ bei einem Doppelklick).
\\ \hline
\texttt{public int getX()} &
x-Position der Maus.
\\ \hline
\texttt{public int getY()} &
y-Position der Maus.
\\ \hline
\end{tabularx}
\egroup


\section{Adapterklassen}

Wenn man ein Interface implementiert, muss jede einzelne Methode des Interfaces
vorhanden sein, auch wenn man vielleicht nur eine einzige Methode benötigt. Das
kann auf die Dauer recht nervig sein. Deshalb stellt die Bibliothek für viele
Interfaces mit mehreren Methoden eine sogenannte \emph{Adapterklasse} bereit.
Eine Adapterklasse ist eine Klasse, die ein Interface mit leeren
Methodenrümpfen implementiert. Zum Beispiel gibt es für das Interface
\myClass{MouseListener} die Klasse \myClass{MouseAdapter}, für das Interface
\myClass{MouseMotionListener} die Klasse \myClass{MouseMotionAdapter} und für
das Interface \myClass{KeyListener} die Klasse \myClass{KeyAdapter}. Alle
Adapterklassen befinden sich im Package \myPackage{java.awt.event}.

Anstatt das Interface vollständig zu implementieren, kann man sich einfach von
der entsprechenden Adapterklasse ableiten und nur die benötigten Methoden
überschreiben. Der Nachteil dabei ist, dass man sich dann von keiner anderen
Klasse mehr ableiten kann, da eine Java-Klasse immer nur eine einzige
Superklasse haben darf (in einigen anderen Programmiersprachen gibt es diese
Regel nicht).
